% GENERATED FILE. Edit canonical lessons, then run npm run generate:books.
% canonical-source-hash: fnv1a64:20d5ae2d6e098b78
% canonical-lessons: JA-C490-ryoushi, JA-C490-gyogyou, JA-C490-gyosen, JA-C490-esa, JA-C490-haneru

\chapter{Fishing Boats}
\label{ch:ja-fishing-boats}
\hlchaptermodality{490}

\begin{chapteropening}
\textbf{By the end of this chapter:} \emph{I can talk about the people who fish for a living.}

Complete the last lesson of chapter 490: I can talk about the people who fish for a living.
\end{chapteropening}

\section[\texorpdfstring{\ja{りょうし} (ryōshi)}{りょうし (ryōshi)}]{\ja{りょうし} (ryōshi) — a fisherman}

\label{lesson:JA-C490-ryoushi}



\begin{quote}
Before the new one: say the Japanese for to ripen, then the Japanese for wheat.
\end{quote}

\subsection*{You'll want to know: \ja{りょうし}}
\textbf{\ja{りょうし}} — \emph{ryōshi} — \textquotedblleft{}a fisherman\textquotedblright{}.

\begin{quote}
\textbf{\ja{りょうしは} \ja{あさ} \ja{はやく} \ja{うみに} \ja{でます。}} — \emph{ryōshi wa asa hayaku umi ni demasu} — \textquotedblleft{}Fishermen go out to sea early in the morning.\textquotedblright{}
\end{quote}

\begin{cousinweb}
\textbf{\ja{りょうし}} (\emph{ryōshi}) is Sino-Japanese (kango): two Chinese characters read with their on'yomi, readings borrowed from Chinese. They are \emph{ryō}, \textquotedblleft{}fishing\textquotedblright{} and \emph{shi}, \textquotedblleft{}expert\textquotedblright{} (Middle Chinese \emph{srij}) (KANJIDIC2; Middle Chinese in Baxter's transcription).
\end{cousinweb}

\begin{grammarlens}[title={What you've built}]
One more word for describing what you lived and felt.
\end{grammarlens}

\subsection*{Your turn}
\begin{itemize}
\raggedright
  \item \emph{Say it:} \emph{ryōshi}
  \item \emph{Say it:} \emph{ryōshi}, once more
  \item \emph{Say it:} \emph{ryōshi wa asa hayaku umi ni demasu} — Fishermen go out to sea early in the morning.
  \item \emph{Recall:} say the Japanese for to ripen, then the Japanese for wheat, then say \emph{ryōshi} again
\end{itemize}

\begin{tcolorbox}[breakable,colback=teal!4,colframe=teal!35!black,title={Before you move on}]
What does \ja{りょうし} mean? (\textquotedblleft{}A fisherman\textquotedblright{}.) Say it once more.
\end{tcolorbox}
\section[\texorpdfstring{\ja{ぎょぎょう} (gyogyō)}{ぎょぎょう (gyogyō)}]{\ja{ぎょぎょう} (gyogyō) — the fishing industry}

\label{lesson:JA-C490-gyogyou}



\begin{quote}
Before the new one: say the Japanese for wheat, then the Japanese for a fisherman.
\end{quote}

\subsection*{You'll want to know: \ja{ぎょぎょう}}
\textbf{\ja{ぎょぎょう}} — \emph{gyogyō} — \textquotedblleft{}the fishing industry\textquotedblright{}.

\begin{quote}
\textbf{\ja{この} \ja{みなとは} \ja{ぎょぎょうの} \ja{まちです。}} — \emph{kono minato wa gyogyō no machi desu} — \textquotedblleft{}This port is a fishing town.\textquotedblright{}
\end{quote}

\begin{cousinweb}
\textbf{\ja{ぎょぎょう}} (\emph{gyogyō}) is Sino-Japanese (kango): two Chinese characters read with their on'yomi, readings borrowed from Chinese. They are \emph{gyo}, \textquotedblleft{}fishing\textquotedblright{} and \emph{gyō}, \textquotedblleft{}business\textquotedblright{} (Middle Chinese \emph{ngjaep}) (KANJIDIC2; Middle Chinese in Baxter's transcription).
\end{cousinweb}

\begin{grammarlens}[title={What you've built}]
One more word for describing what you lived and felt.
\end{grammarlens}

\subsection*{Your turn}
\begin{itemize}
\raggedright
  \item \emph{Say it:} \emph{gyogyō}
  \item \emph{Say it:} \emph{gyogyō}, once more
  \item \emph{Say it:} \emph{kono minato wa gyogyō no machi desu} — This port is a fishing town.
  \item \emph{Recall:} say the Japanese for wheat, then the Japanese for a fisherman, then say \emph{gyogyō} again
\end{itemize}

\begin{tcolorbox}[breakable,colback=teal!4,colframe=teal!35!black,title={Before you move on}]
What does \ja{ぎょぎょう} mean? (\textquotedblleft{}The fishing industry\textquotedblright{}.) Say it once more.
\end{tcolorbox}
\section[\texorpdfstring{\ja{ぎょせん} (gyosen)}{ぎょせん (gyosen)}]{\ja{ぎょせん} (gyosen) — a fishing boat}

\label{lesson:JA-C490-gyosen}



\begin{quote}
Before the new one: say the Japanese for a fisherman, then the Japanese for the fishing industry.
\end{quote}

\subsection*{You'll want to know: \ja{ぎょせん}}
\textbf{\ja{ぎょせん}} — \emph{gyosen} — \textquotedblleft{}a fishing boat\textquotedblright{}.

\begin{quote}
\textbf{\ja{みなとに} \ja{ぎょせんが} \ja{ならんでいます。}} — \emph{minato ni gyosen ga narandeimasu} — \textquotedblleft{}Fishing boats are lined up in the harbour.\textquotedblright{}
\end{quote}

\begin{cousinweb}
\textbf{\ja{ぎょせん}} (\emph{gyosen}) is Sino-Japanese (kango): two Chinese characters read with their on'yomi, readings borrowed from Chinese. They are \emph{gyo}, \textquotedblleft{}fishing\textquotedblright{} and \emph{sen}, \textquotedblleft{}ship\textquotedblright{} (Middle Chinese \emph{zywen}) (KANJIDIC2; Middle Chinese in Baxter's transcription).
\end{cousinweb}

\begin{grammarlens}[title={What you've built}]
One more word for describing what you lived and felt.
\end{grammarlens}

\subsection*{Your turn}
\begin{itemize}
\raggedright
  \item \emph{Say it:} \emph{gyosen}
  \item \emph{Say it:} \emph{gyosen}, once more
  \item \emph{Say it:} \emph{minato ni gyosen ga narandeimasu} — Fishing boats are lined up in the harbour.
  \item \emph{Recall:} say the Japanese for a fisherman, then the Japanese for the fishing industry, then say \emph{gyosen} again
\end{itemize}

\begin{tcolorbox}[breakable,colback=teal!4,colframe=teal!35!black,title={Before you move on}]
What does \ja{ぎょせん} mean? (\textquotedblleft{}A fishing boat\textquotedblright{}.) Say it once more.
\end{tcolorbox}
\section[\texorpdfstring{\ja{えさ} (esa)}{えさ (esa)}]{\ja{えさ} (esa) — animal feed; bait}

\label{lesson:JA-C490-esa}



\begin{quote}
Before the new one: say the Japanese for the fishing industry, then the Japanese for a fishing boat.
\end{quote}

\subsection*{You'll want to know: \ja{えさ}}
\textbf{\ja{えさ}} — \emph{esa} — \textquotedblleft{}animal feed; bait\textquotedblright{}.

\begin{quote}
\textbf{\ja{ねこの} \ja{えさを} \ja{かいます。}} — \emph{neko no esa o kaimasu} — \textquotedblleft{}I buy food for the cat.\textquotedblright{}
\end{quote}

\begin{cousinweb}
\textbf{\ja{えさ}} (\emph{esa}) is a native Japanese word (wago), inherited rather than borrowed from Chinese; its kanji was attached later as a kun'yomi, and its older origin is not settled (JMdict, KANJIDIC2).
\end{cousinweb}

\begin{grammarlens}[title={What you've built}]
One more word for describing what you lived and felt.
\end{grammarlens}

\subsection*{Your turn}
\begin{itemize}
\raggedright
  \item \emph{Say it:} \emph{esa}
  \item \emph{Say it:} \emph{esa}, once more
  \item \emph{Say it:} \emph{neko no esa o kaimasu} — I buy food for the cat.
  \item \emph{Recall:} say the Japanese for the fishing industry, then the Japanese for a fishing boat, then say \emph{esa} again
\end{itemize}

\begin{tcolorbox}[breakable,colback=teal!4,colframe=teal!35!black,title={Before you move on}]
What does \ja{えさ} mean? (\textquotedblleft{}Animal feed; bait\textquotedblright{}.) Say it once more.
\end{tcolorbox}
\section[\texorpdfstring{\ja{はねる} (haneru)}{はねる (haneru)}]{\ja{はねる} (haneru) — to jump; to leap}

\label{lesson:JA-C490-haneru}



\begin{quote}
Before the new one: say the Japanese for a fishing boat, then the Japanese for animal feed.
\end{quote}

\subsection*{You'll want to know: \ja{はねる}}
\textbf{\ja{はねる}} — \emph{haneru} — \textquotedblleft{}to jump; to leap\textquotedblright{}.

\begin{quote}
\textbf{\ja{さかなが} \ja{みずの} \ja{なかで} \ja{はねました。}} — \emph{sakana ga mizu no naka de hanemashita} — \textquotedblleft{}A fish jumped in the water.\textquotedblright{}
\end{quote}

Polite \textbf{\ja{はねます}}, \emph{hanemasu}; an ichidan verb. Past (polite): \textbf{\ja{さかなが} \ja{はねました。}}, \emph{sakana ga hanemashita}.

\begin{cousinweb}
\textbf{\ja{はねる}} (\emph{haneru}) is a native Japanese word (wago), inherited rather than borrowed from Chinese, and usually written in kana; its older origin is not settled (JMdict).
\end{cousinweb}

\begin{grammarlens}[title={What you've built}]
One more word for describing what you lived and felt.
\end{grammarlens}

\subsection*{Your turn}
\begin{itemize}
\raggedright
  \item \emph{Say it:} \emph{haneru}
  \item \emph{Say it:} \emph{haneru}, once more
  \item \emph{Say it:} \emph{sakana ga mizu no naka de hanemashita} — A fish jumped in the water.
  \item \emph{Recall:} say the Japanese for a fishing boat, then the Japanese for animal feed, then say \emph{haneru} again
\end{itemize}

\begin{tcolorbox}[breakable,colback=teal!4,colframe=teal!35!black,title={Before you move on}]
What does \ja{はねる} mean? (\textquotedblleft{}To jump; to leap\textquotedblright{}.) Say it once more.
\end{tcolorbox}
